\begin{helper}[The elementary inequality after (4.4)]
Let $p>1$.  Then there is a constant $C=C(p)>0$ such that for every integer
$n\ge1$, every $\varepsilon\ge0$ and all $a,b\in\R^n$,
\[
\int_0^1(1-t)\bigl(\varepsilon+|ta+(1-t)b|^2\bigr)^{\frac p2-1}dt
\;\ge\;C\,\bigl(\varepsilon+|a|^2+|b|^2\bigr)^{\frac p2-1} .
\]
All powers are real powers of nonnegative reals.

The constant is chosen before the dimension, which is what the paper's
``where $C(p)$ is a constant depending only on $p$'' (thesis p.~34, the
sentence following (4.4) of the thesis) asserts: $C$ may depend on $p$ and on nothing
else, $n$ included.  The proof below produces such a $C$ explicitly.
\end{helper}
\begin{proof}
Write $s=\frac p2-1$, $\varphi(t)=ta+(1-t)b$ for $t\in[0,1]$, and
$T=\varepsilon+|a|^2+|b|^2$.  Note $\varphi(0)=b$, $\varphi(1)=a$ and
$\varphi(t)-\varphi(t')=(t-t')(a-b)$.  Set $M=\max\{|a|,|b|\}$ and $L=|a-b|$.

\emph{Why the constant may be chosen before the dimension.}  Every quantity
occurring in the two sides of the asserted inequality is a function of the three
real numbers $|a|$, $|b|$ and $(a,b)$ alone, with no other dependence on $n$ or
on $a,b$ themselves.  Indeed, expanding the inner product,
\[
|\varphi(t)|^2=|t\,a+(1-t)b|^2
=t^2|a|^2+2t(1-t)\,(a,b)+(1-t)^2|b|^2
\qquad(t\in\R),
\]
which uses only $|x+y|^2=|x|^2+2(x,y)+|y|^2$, the bilinearity of the inner
product and its symmetry; and $T=\varepsilon+|a|^2+|b|^2$ is already of that
form.  So for a given triple $(|a|,|b|,(a,b))$ the assertion is literally the
same assertion in every dimension $n\ge1$ that realises that triple, and a
constant proved to work for all such triples works in every dimension at once.
The constants produced in the three cases below are built only from $p$ and the
numerical factors $2$, $\tfrac12$, $\tfrac34$, $\tfrac14$, $8$, $64$, $128$ and
$16$, never from $n$, which is the concrete form of that observation; the
estimates themselves use only the triangle inequality $|x+y|\le|x|+|y|$ and
$|a|^2+|b|^2\le2M^2$, both valid in every dimension.

\emph{Degenerate case.}  If $T=0$ then $\varepsilon=0$ and $a=b=0$, so
$\varphi\equiv0$ and both sides equal $C\cdot0^{s}$ resp.\
$\int_0^1(1-t)0^{s}dt$, which are $0$ when $s\ne0$ and $C$ resp.\ $\tfrac12$
when $s=0$.  The inequality therefore holds in this case as soon as
$C\le\tfrac12$, which the choices below respect.  Assume from now on $T>0$.

\emph{Case 1: $1<p<2$, i.e.\ $s<0$.}  For every $t\in[0,1]$,
\[
|\varphi(t)|\le t|a|+(1-t)|b|\le|a|+|b| ,
\]
hence $|\varphi(t)|^2\le(|a|+|b|)^2\le2(|a|^2+|b|^2)$ and therefore
\[
0<\varepsilon+|\varphi(t)|^2\le\varepsilon+2(|a|^2+|b|^2)\le 2T .
\]
(The left strict inequality holds because if $\varepsilon+|\varphi(t)|^2=0$ then
$\varepsilon=0$ and $\varphi(t)=0$; since $s<0$ we may in that case replace the
integrand by $+\infty$ at that single point, so the bound below is only
improved. More precisely, the function $t\mapsto\varepsilon+|\varphi(t)|^2$
vanishes for at most one $t$, and for all other $t$ the inequality
$(\varepsilon+|\varphi(t)|^2)^{s}\ge(2T)^{s}$ holds; a single point does not
affect the integral.)  Since $s<0$, raising to the power $s$ reverses the
inequality:
\[
\bigl(\varepsilon+|\varphi(t)|^2\bigr)^{s}\ge(2T)^{s}=2^{s}T^{s}
\qquad\text{for a.e.\ }t\in[0,1].
\]
Before integrating we record that the integrand is interval-integrable on
$[0,1]$, so that the monotonicity of the integral may be applied to it: by
\noderef{SegmentWeightIntegrable}, with the present $n,p,\varepsilon,a,b$ --- its
hypotheses are exactly $p>1$ and $\varepsilon\ge0$, both of which we have --- the
weight $t\mapsto(\varepsilon+|\varphi(t)|^2)^{s}$ is interval-integrable with
respect to Lebesgue measure on $[0,1]$, and multiplying it by the continuous
bounded factor $1-t$ preserves interval-integrability.  This is the one place in
the proof where finiteness is not automatic: in the present case $s<0$ and the
weight is unbounded near a zero of $\varepsilon+|\varphi|^2$, so without
\noderef{SegmentWeightIntegrable} the left-hand side of the asserted inequality
need not be a well-defined finite number, and comparing it with
$\tfrac12 2^{s}T^{s}$ would compare a junk value.  (In Case 2 below $s\ge0$ and
the integrand is continuous on $[0,1]$, hence integrable with no further
argument.)

Multiplying by $(1-t)\ge0$ and integrating, and using
$\int_0^1(1-t)\,dt=\tfrac12$,
\[
\int_0^1(1-t)\bigl(\varepsilon+|\varphi(t)|^2\bigr)^{s}dt\ \ge\ \tfrac12\,2^{s}\,T^{s} .
\]
So the claim holds in this case with $C=2^{s-1}=2^{p/2-2}$.

\emph{Case 2: $p\ge2$, i.e.\ $s\ge0$.}  We first produce a set of $t$'s of
substantial measure on which $|\varphi(t)|$ is bounded below by a fixed multiple
of $M$.

\emph{Subcase 2a: $L<M/2$.}  Suppose $|a|=M$ (the case $|b|=M$ is symmetric).
For every $t\in[0,1]$, $|\varphi(t)-a|=(1-t)L\le L<M/2$, hence
\[
|\varphi(t)|\ \ge\ |a|-|\varphi(t)-a|\ >\ M-\tfrac M2=\tfrac M2 .
\]
Thus $|\varphi(t)|\ge M/2$ for every $t\in[0,1]$, and we may take
$E=[0,\tfrac34]$, which has measure $\tfrac34\ge\tfrac14$, with
$|\varphi(t)|\ge M/8$ on $E$ (indeed $M/2\ge M/8$).

\emph{Subcase 2b: $L\ge M/2$.}  Let
$Z=\{t\in[0,1]:|\varphi(t)|\le M/8\}$.  If $t,t'\in Z$ then
\[
|t-t'|\,L=|\varphi(t)-\varphi(t')|\le|\varphi(t)|+|\varphi(t')|\le\tfrac M4 ,
\]
so $|t-t'|\le M/(4L)\le M/(4\cdot M/2)=\tfrac12$ (if $M=0$ then $a=b=0$,
$L=0$, and $T=\varepsilon$; this degenerate situation is covered by Subcase 2a
with the bound $|\varphi(t)|\ge0$ and is treated at the end of this case).
Hence $Z$ is contained in an interval of length $\tfrac12$, so its Lebesgue
measure is at most $\tfrac12$.  Put $E=[0,\tfrac34]\setminus Z$; then
$|E|\ge\tfrac34-\tfrac12=\tfrac14$ and $|\varphi(t)|>M/8$ on $E$.

In both subcases we have a measurable $E\subseteq[0,\tfrac34]$ with
$|E|\ge\tfrac14$ and $|\varphi(t)|\ge M/8$ for $t\in E$.  On $E$,
\[
\varepsilon+|\varphi(t)|^2\ \ge\ \varepsilon+\tfrac{M^2}{64}
\ \ge\ \tfrac1{64}\bigl(\varepsilon+M^2\bigr)
\ \ge\ \tfrac1{128}\bigl(\varepsilon+|a|^2+|b|^2\bigr)=\tfrac{T}{128},
\]
where the last step uses $|a|^2+|b|^2\le2M^2$, so
$\varepsilon+|a|^2+|b|^2\le\varepsilon+2M^2\le2(\varepsilon+M^2)$.  Since
$s\ge0$, raising to the power $s$ preserves the inequality:
$(\varepsilon+|\varphi(t)|^2)^{s}\ge128^{-s}T^{s}$ on $E$.  Moreover
$1-t\ge\tfrac14$ on $E\subseteq[0,\tfrac34]$.  Since the integrand is
nonnegative (both factors are $\ge0$), restricting the integral to $E$ gives
\[
\int_0^1(1-t)\bigl(\varepsilon+|\varphi(t)|^2\bigr)^{s}dt
\ \ge\ \int_E\tfrac14\cdot128^{-s}T^{s}\,dt
\ \ge\ \tfrac14\cdot\tfrac14\cdot128^{-s}\,T^{s}
=\tfrac1{16}\,128^{-s}\,T^{s} .
\]
If $M=0$ then $a=b=0$, $\varphi\equiv0$, $T=\varepsilon>0$, and
$\varepsilon+|\varphi(t)|^2=\varepsilon=T$ for every $t$, so the integral equals
$\tfrac12T^{s}$, which is at least $\tfrac1{16}128^{-s}T^{s}$.  So the claim
holds in Case 2 with $C=\tfrac1{16}128^{-s}=\tfrac1{16}128^{1-p/2}$.

Taking $C=\min\{2^{p/2-2},\ \tfrac1{16}128^{1-p/2},\ \tfrac12\}>0$, which depends
only on $p$ and in particular not on $n$, $\varepsilon$, $a$ or $b$, proves the
lemma in all cases; the hypothesis $n\ge1$ was used nowhere beyond guaranteeing
that $\R^n$ is a nonzero space, and the displayed choice of $C$ is made once,
before any dimension is named.
\end{proof}
